\documentclass[aspectratio=169, xcolor=table]{beamer}

% Professional Theme
\usetheme{Madrid}
\usecolortheme{whale}

% Packages
\usepackage{graphicx}
\usepackage{booktabs}
\usepackage{tikz}
\usepackage{hyperref}

% Title Page Information
\title[ML in Marine Geophysics \& Oceanography]{Predictive Machine Learning in Marine Geophysics and Oceanography}
\subtitle{AI Solutions for Gravity Anomaly Resolution and ENSO-Aware Current Forecasting}
\author[Shivanshu Tiwari]{Shivanshu Tiwari \\ \vspace{0.5cm} \footnotesize Supervisor: Pranjeet Sarkar}
\date{September 2026}

\begin{document}

% Slide 1: Title
\begin{frame}
    \titlepage
\end{frame}

% Slide 2: Objectives
\begin{frame}{Project Objectives}
    \textbf{Core Goal:} To develop two state-of-the-art physics-informed deep learning models to solve critical resolution and forecasting gaps in oceanography.
    \vspace{0.5cm}
    \begin{enumerate}
        \item \textbf{Geophysics:} Create a Super-Resolution ML model to predict high-resolution marine gravity anomalies from low-resolution datasets.
        \vspace{0.2cm}
        \item \textbf{Oceanography:} Build a Spatiotemporal Forecasting model to predict ocean currents while dynamically accounting for complex climate events like El Ni\~no (ENSO).
    \end{enumerate}
\end{frame}

% Slide 3: Problem 1
\begin{frame}{Part 1: The Gravity Anomaly Problem}
    \begin{columns}
        \begin{column}{0.5\textwidth}
            \textbf{The Challenge:}
            \begin{itemize}
                \item Low-resolution Bouguer-like gravity maps are readily available globally.
                \item High-resolution maps (crucial for identifying hidden seamounts and trenches) are difficult and expensive to acquire via shipborne surveys.
            \end{itemize}
        \end{column}
        \begin{column}{0.5\textwidth}
            \centering
            % Placeholder for the Global Map
            \includegraphics[width=0.9\textwidth, height=0.6\textheight, keepaspectratio]{C:/Users/shiva/Downloads/Indian_Ocean_Features/grav_SWOT_05_Global.png} \\
            \footnotesize \textit{Current Global High-Res Coverage}
        \end{column}
    \end{columns}
\end{frame}

% Slide 4: Solution 1
\begin{frame}{Part 1: The Machine Learning Solution}
    \begin{columns}
        \begin{column}{0.5\textwidth}
            \textbf{Physics-Informed Super Resolution:}
            \begin{itemize}
                \item \textbf{Method:} Use easily available low-res maps as input to predict uncharted high-res anomalies.
                \item \textbf{Ground Truth:} Train on cutting-edge SWOT V05 data (1 arc-minute resolution).
                \item \textbf{Outcome:} Extrapolate precise geological features globally.
            \end{itemize}
        \end{column}
        \begin{column}{0.5\textwidth}
            \centering
            \includegraphics[width=0.9\textwidth, height=0.6\textheight, keepaspectratio]{C:/Users/shiva/Downloads/Indian_Ocean_Features/grav_SWOT_05_IndianOcean_Features_Colored.png} \\
            \footnotesize \textit{High-Res Features (Seamounts/Trenches)}
        \end{column}
    \end{columns}
\end{frame}

% Slide 5: Architecture 1
\begin{frame}{Part 1: Architecture (Swin Transformer)}
    \textbf{State-of-the-Art Design (2026): Physics-Informed Swin Transformer}
    \vspace{0.3cm}
    \begin{itemize}
        \item \textbf{Input:} Low-res Gravity, Bathymetry, Vertical Deflections (4 channels).
        \item \textbf{Core:} 6 Residual Swin Transformer Blocks (RSTB) to capture global spatial dependencies via Self-Attention.
        \item \textbf{Physics Constraint:} Custom Spectral Loss Layer enforcing Laplace's equation ($\nabla^2 g = 0$). Prevents physically impossible predictions.
        \item \textbf{Output:} 1 arc-minute High-Resolution Gravity map.
    \end{itemize}
\end{frame}

% Slide 6: Problem 2
\begin{frame}{Part 2: The Ocean Current Problem}
    \begin{columns}
        \begin{column}{0.5\textwidth}
            \textbf{The Challenge:}
            \begin{itemize}
                \item Ocean currents are generally predictable, but weather disturbances (El Ni\~no/ENSO) cause massive, non-linear variations in velocity and temperature.
                \item Traditional physics-based models (NEMO) struggle to adapt to these rapid climate-driven anomalies.
            \end{itemize}
        \end{column}
        \begin{column}{0.5\textwidth}
            \centering
            \includegraphics[width=0.9\textwidth, height=0.6\textheight, keepaspectratio]{C:/Users/shiva/Downloads/Indian_Ocean_Features/Ocean_Current_Map.png} \\
            \footnotesize \textit{Surface Current Velocities}
        \end{column}
    \end{columns}
\end{frame}

% Slide 7: Solution 2
\begin{frame}{Part 2: The Machine Learning Solution}
    \textbf{Spatiotemporal Forecasting with Climate Embedding}
    \vspace{0.3cm}
    \begin{itemize}
        \item \textbf{Method:} Train on a massive 5.5-year 4D dataset (Copernicus CMEMS) encompassing historical ENSO cycles.
        \item \textbf{Innovation:} Explicitly embed real-time ENSO (Ni\~no 3.4) and IOD indices directly into the neural network to teach the model "teleconnections".
        \item \textbf{Outcome:} Highly accurate predictions of $u$ and $v$ current velocities at multiple depths for $T+1, T+3, T+6$ months.
    \end{itemize}
\end{frame}

% Slide 8: Architecture 2
\begin{frame}{Part 2: Architecture (PhyGOSCR-Net)}
    \textbf{Physics-Guided Ocean Surface Current Reconstruction Network}
    \vspace{0.2cm}
    \begin{itemize}
        \item \textbf{Spatial Encoder:} Multi-scale Dilated CNN to capture local eddies (dilation=1), mesoscale gyres (dilation=3), and basin-wide ENSO signals (dilation=9).
        \item \textbf{Temporal Encoder:} 4-Layer Transformer Encoder to process the 12-month rolling history.
        \item \textbf{Physics Constraints:} 
        \begin{itemize}
            \item \textit{Geostrophic Balance Loss:} Enforces physical relationship between Sea Surface Height and Current Velocity.
            \item \textit{Non-Divergence Loss:} Enforces fluid incompressibility.
        \end{itemize}
    \end{itemize}
\end{frame}

% Slide 9: Data Readiness
\begin{frame}{Data Readiness \& Acquisition}
    \textbf{100\% of required datasets have been successfully acquired and stored locally:}
    \vspace{0.3cm}
    \begin{table}[]
        \centering
        \resizebox{\textwidth}{!}{
        \begin{tabular}{@{}llc@{}}
        \toprule
        \textbf{Dataset} & \textbf{Source} & \textbf{Status} \\ \midrule
        High-Res Gravity (SWOT V05) & UCSD / Sandwell & \textbf{Ready} \\
        3D Ocean Currents ($u, v$) 2020-2026 & Copernicus CMEMS & \textbf{Ready} (40 GB) \\
        Ocean Temperature \& SSH & Copernicus CMEMS & \textbf{Ready} \\
        Ocean Salinity & Copernicus CMEMS & \textbf{Ready} (15 GB) \\
        Surface Wind Stress & ECMWF ERA5 & \textbf{Ready} \\
        ENSO \& IOD Climate Indices & NOAA / JAMSTEC & \textbf{Ready} \\ \bottomrule
        \end{tabular}
        }
    \end{table}
\end{frame}

% Slide 10: Conclusion
\begin{frame}{Summary \& Next Steps}
    \textbf{Summary:}
    \begin{itemize}
        \item We have architected two state-of-the-art, physics-informed AI models.
        \item The models directly address the supervisor's requirements regarding high-resolution gravity mapping and ENSO-aware current forecasting.
        \item All necessary data (over 60 GBs of high-fidelity geospatial and oceanographic records) is fully downloaded and prepared.
    \end{itemize}
    \vspace{0.5cm}
    \textbf{Next Steps (Post-Approval):}
    \begin{itemize}
        \item Implement PyTorch architectures.
        \item Execute data fusion and preprocessing pipelines.
        \item Commence GPU training and validation.
    \end{itemize}
\end{frame}

\end{document}
